

O Brasil, segundo o Instituto Brasileiro de Geografia e Estatística (IBGE), em 2014, a estimativa de depressão entre indivíduos entre 18 a 29 anos corresponde a 3,9\% da população, correspondendo a um grupo de jovens adultos. Boa parte disso se deve à necessidade de tomada de decisões, pressão social, entre outros fatores \cite{souza_condicoes_2017}.

Esse período é conhecido por possuir um dos maiores riscos psicológicos, podendo gerar gatilhos emocionais e estresse duradouro \cite{silveira_saumental_2011}. Também citado por \citeonline{bolsoni-silva_o_2014}, a universidade é capaz de agravar problemas de saúde pré-existentes ou aumentar a probabilidade do seu surgimento. Sendo assim, a universidade é um ambiente estressante, principalmente no primeiro semestre. A solidão, um fator crucial para depressão, ocorre principalmente no primeiro ano da universidade devido a um novo ambiente e à falta de experiências anteriores em contextos desafiadores e exigentes \cite{rhodes_loneliness_2014}.

Ressalta-se a importância dos hábitos saudáveis, lazer, presença religiosa e atividade física, que ajudam a promover o bem-estar e são possíveis formas de lidar com estados emocionais negativos. Outro fator crucial para o desenvolvimento de problemas psicológicos é a falta de hábitos saudáveis e menor qualidade de vida nos domínios psicológico e ambiental \cite{barroso_solidao_2019}.

Além disso, \citeonline{barroso_solidao_2019} destaca como os hábitos de vida influenciam diretamente a relação com a solidão e a depressão. Comportamentos como acordar cansado e realizar menos refeições ao longo do dia, entre outros, revelaram um impacto significativo na qualidade de vida dos participantes.

Outro aspecto que influencia a saúde emocional é o suporte social que, como defendido por \citeonline{siqueira_construcao_2008} em "Construção e validação da escala de percepção de suporte social", é definido como “o processo no qual alguém pode conseguir ajuda de ordem instrumental, financeira ou emocional”. Esse apoio é considerado essencial para a promoção e manutenção da saúde mental.